\begin{figure}[H]
    \centering
    \begin{tikzpicture}[x=1.28cm, y=1.0cm]

        % (a) naive transfer: one global pass per layer
        \node[panel] at (-1.75,1.30) {(α) άμεση μεταφορά};

        \foreach \i/\lab in {0/{$2^{20}$},1/{$2^{19}$},2/{$\cdots$},3/{$32$},4/{$16$},5/{$8$},6/{$4$},7/{$2$},8/{$1$}} {
            \node[lane, minimum width=1.02cm] (a\i) at (\i,0.35) {\lab};
        }
        \node[rowlab] at (a0.west) {επίπεδο ($j$)};

        \node[blkdrm, minimum width=11.7cm, minimum height=0.46cm] (adram) at (4,-1.10) {κύρια μνήμη};
        \foreach \i in {0,...,8} {
            \draw[flowmute] ([xshift=-3pt]a\i.south) -- ([xshift=-3pt]a\i.south|-adram.north);
            \draw[flowmute] ([xshift=3pt]a\i.south|-adram.north) -- ([xshift=3pt]a\i.south);
        }
        \node[note, anchor=north] at (4,-1.48) {εννέα περάσματα, $2n$ το καθένα};

        \begin{scope}[yshift=-4.05cm]

            % (b) fused: layers grouped by distance j into a single pass
            \node[panel] at (-1.75,1.62) {(β) συγχώνευση επιπέδων};

            \foreach \i/\lab in {0/{$2^{20}$},1/{$2^{19}$},2/{$\cdots$}} {
                \node[lanedrm, minimum width=1.02cm] (b\i) at (\i,0.35) {\lab};
            }
            \node[lanecac, minimum width=1.02cm] (b3) at (3,0.35) {$32$};
            \foreach \i/\lab in {4/{$16$},5/{$8$},6/{$4$},7/{$2$},8/{$1$}} {
                \node[lanereg, minimum width=1.02cm] (b\i) at (\i,0.35) {\lab};
            }
            \node[rowlab] at (b0.west) {επίπεδο ($j$)};

            \draw[brace] (-0.44,0.66) -- (2.44,0.66);
            \draw[brace] (2.56,0.66) -- (3.44,0.66);
            \draw[brace] (3.56,0.66) -- (8.44,0.66);
            \node[notemute, font=\tiny] at (1.0,1.02) {$j>$ τμήμα};
            \node[notemute, font=\tiny] at (3.0,1.02) {$j\le$ τμήμα};
            \node[notemute, font=\tiny] at (6.0,1.02) {$j<W$};

            \node[blkdrm, minimum width=11.7cm, minimum height=0.46cm] (bdram) at (4,-1.10) {κύρια μνήμη};

            \foreach \i in {0,1,2} {
                \draw[flowmute] ([xshift=-3pt]b\i.south) -- ([xshift=-3pt]b\i.south|-bdram.north);
                \draw[flowmute] ([xshift=3pt]b\i.south|-bdram.north) -- ([xshift=3pt]b\i.south);
            }
            \draw[flowmute] ([xshift=-3pt]b3.south) -- ([xshift=-3pt]b3.south|-bdram.north);
            \draw[flowmute] ([xshift=3pt]b3.south|-bdram.north) -- ([xshift=3pt]b3.south);
            \draw[flowmute] ([xshift=-3pt]b6.south) -- ([xshift=-3pt]b6.south|-bdram.north);
            \draw[flowmute] ([xshift=3pt]b6.south|-bdram.north) -- ([xshift=3pt]b6.south);

            \node[blkcac, minimum width=1.15cm, minimum height=0.40cm] at (3,-0.42) {\tiny τμήμα};
            \node[blkreg, minimum width=1.35cm, minimum height=0.40cm] at (6,-0.42) {\tiny καταχ.};

            \node[note, anchor=north] at (4,-1.48) {πέντε περάσματα};
        \end{scope}

    \end{tikzpicture}
    \caption[Συγχώνευση διαδοχικών επιπέδων σε ένα πέρασμα]{Η συγχώνευση διαδοχικών επιπέδων σε ένα πέρασμα· κάθε ζεύγος βελών
    είναι μία ανάγνωση και μία εγγραφή ολόκληρου του πίνακα. Στην άμεση μεταφορά
    (α) κάθε επίπεδο του χρονοδιαγράμματος είναι ξεχωριστό πέρασμα, με ελάχιστο
    υπολογισμό ανάμεσα σε δύο διαδοχικές διαδρομές προς τη μνήμη. Στη συγχωνευμένη
    υλοποίηση (β) η απόσταση $j$ καθορίζει πού μπορεί να εκτελεστεί το επίπεδο: για
    $j$ μικρότερο του πλάτους $W$ του καταχωρητή τα δύο στοιχεία κάθε σύγκρισης
    βρίσκονται ήδη στον ίδιο καταχωρητή και ολόκληρη η ακολουθία εκτελείται ως μία
    φόρτωση και μία αποθήκευση· για $j$ έως το μέγεθος του τμήματος η ακολουθία
    εκτελείται τμήμα προς τμήμα, με το τμήμα να παραμένει στην κρυφή μνήμη για όλα
    τα επίπεδα· μόνο τα επίπεδα με πολύ μεγάλο $j$ διατηρούν ξεχωριστό πέρασμα. Οι
    αριθμοί του σχήματος είναι ενδεικτικοί.}
    \label{fig:cpu-layer-fusion}
\end{figure}
